% !TeX root = Chapter_02_Slides.tex
% Chapter 2 lecture slides — Enterprise Risk Management: Evolution, Frameworks, and Institutional Context
% Instructor-resource series.
\documentclass[11pt, aspectratio=169]{beamer}
\usepackage{tikz}
\makeatletter
\def\input@path{{./}{../slides/}}
\makeatother
\usepackage{erm_beamer_theme}
\usepackage{amsmath,amssymb}

\title[ERM Ch.\,2]{Enterprise Risk Management: Evolution, Frameworks, and Institutional Context}
\subtitle{Enterprise Risk Management, Second Edition --- Chapter 2}
\author{Martin F.\ Grace}
\institute{University of Iowa\\[2pt] Tippie College of Business\\[2pt] Vaughan Institute of Risk Management}
\date{June 2026}

\begin{document}

\begin{frame}[plain]
\titlepage
\end{frame}

\begin{frame}{Learning Objectives}
\begin{itemize}
  \item Trace the arc of corporate risk management from \alert{insurance purchasing} to enterprise-wide strategy.
  \item Explain how Sarbanes-Oxley and major corporate failures created the impetus for ERM.
  \item Describe the structure and core logic of the \alert{2017 COSO ERM Framework}.
  \item Compare COSO and ISO~31000 --- origin, emphasis, practical application.
  \item Identify the institutions that shape ERM practice in financial and non-financial sectors.
  \item Outline a realistic phased approach to ERM implementation.
\end{itemize}
\end{frame}

\section{Opening Case: The GM Ignition Switch Crisis}

\begin{frame}{It Began with a Small Spring}
\begin{itemize}
  \item A coiled spring in the Chevrolet Cobalt's ignition switch was supposed to hold the key in ``run'' --- engineers knew \emph{before launch} that it was too weak.
  \item A heavy keychain, a knee, a rough road: the switch slips, the engine cuts out, power steering goes --- and the \alert{airbags don't deploy}.
  \item GM first documented the problem in \textbf{2001}. The Cobalt went on sale in 2004.
  \item The recall came in \textbf{February 2014} --- thirteen years later.
\end{itemize}
\end{frame}

\begin{frame}{Thirteen Years of Not Connecting the Dots}
\begin{itemize}
  \item Engineers analyzed stalling complaints. Service bulletins went to dealers. Customers reported cars dying on the highway.
  \item Lawyers settled wrongful-death claims where airbags inexplicably failed. Safety committees filed reports with NHTSA.
  \item Each group worked its lane --- and saw only the slice that landed on its desk.
  \item The question \emph{``Is this a single enterprise-level safety risk requiring a decision at the top?''} belonged to \alert{no one}. So it went unasked.
\end{itemize}
\medskip
The toll: 100+ deaths linked to the defect, a \$900M DOJ settlement, a \$2.5B victim compensation fund, congressional hearings.
\end{frame}

\begin{frame}{The Valukas Verdict}
\begin{itemize}
  \item The independent investigation (Valukas Report, 2014) found no single villain who knew everything and stayed silent.
  \item Instead: many competent people each managing a \emph{piece} of a serious risk, no one responsible for the \emph{whole} --- and the institution drifted for a decade while people died.
  \item Not a shortage of smart people or activity. A shortage of \alert{integration}.
\end{itemize}
\medskip
\begin{block}{The structural question}
The information existed. The expertise existed. What structure would have forced the scattered pieces into one picture --- and who at GM should have owned the enterprise-level question?
\end{block}
\end{frame}

\section{The Evolution of Corporate Risk Management}

\begin{frame}{From Insurance Buying to Integration}
\begin{itemize}
  \item Most of the 20th century: risk management = \alert{insurance management}. Identify insurable exposures, buy coverage.
  \item Risk treated as a cost to minimize, not a strategic consideration --- the function sat well below the executive suite.
  \item 1970s--80s: captives and self-insurance --- financial innovations that cut the cost of risk transfer, but strategy still happened in a different room.
  \item ERM is the response to a recurring failure mode: risks managed in \alert{silos}, with no one integrating the enterprise-level view.
\end{itemize}
\end{frame}

\begin{frame}{The 1990s: Object Lessons in Siloed Finance}
\begin{itemize}
  \item \textbf{Barings Bank} (1995): one trader in Singapore, $\sim$\$1.4B in unauthorized derivatives positions nobody at the top understood or monitored.
  \item \textbf{Orange County} (1994): \$1.6B+ lost on leveraged interest-rate bets the treasurer wasn't authorized to make.
  \item \textbf{Metallgesellschaft} (1993): long-term supply contracts vs.\ short-term hedges --- a cash crisis the supervisory board never saw coming.
\end{itemize}
\medskip
\begin{block}{The common thread}
None of these risks was exotic or unforeseeable in principle. They were \alert{invisible to the people with authority to act} because no one was assigned to look across the whole picture.
\end{block}
\end{frame}

\begin{frame}{The Early 2000s: Governance Scandals}
\begin{itemize}
  \item \textbf{Enron} (2001): accounting fraud, off-balance-sheet structures concealing exposure, a culture that rewarded aggressive engineering and suppressed dissent.
  \item \textbf{WorldCom} (2002): $\sim$\$11B in accounting fraud.
  \item Both had internal audit functions, board audit committees, and external auditors --- \emph{none of it stopped the bleeding}.
  \item The governance structures were inadequate to the risks the companies were actually running.
\end{itemize}
\end{frame}

\section{The Sarbanes-Oxley Inflection}

\begin{frame}{SOX (2002): The Two Provisions That Matter}
\begin{itemize}
  \item \textbf{Section 302}: CEO and CFO personally certify, every quarter, that financial statements fairly present the company's condition and that they own internal controls. \alert{Criminal penalties} for false certification --- executive ignorance became legally dangerous.
  \item \textbf{Section 404}: management must assess internal control effectiveness annually; external auditors must attest. Documentation, control testing, remediation.
  \item SOX did \emph{not} mandate ERM --- but it built the infrastructure (executive accountability, control assessment, board oversight) that made ERM feasible and expected.
\end{itemize}
\end{frame}

\begin{frame}{From Disclosure to Prescription}
\begin{itemize}
  \item Pre-SOX philosophy: require disclosure, trust market discipline to reward good governance and punish bad.
  \item 2001--02 showed disclosure fails when information is too complex to parse, auditors are conflicted, and boards oversee in name only.
  \item SOX moved toward \alert{prescription}: define what controls must look like, hold executives personally responsible, subject auditors to independent oversight (PCAOB).
\end{itemize}
\medskip
\begin{block}{The spillover}
If boards must rigorously oversee financial reporting controls, it's a short step to asking what \emph{other} significant risks are being managed --- and how. COSO's ERM framework (2004) arrived at exactly that moment.
\end{block}
\end{frame}

\section{The COSO ERM Framework}

\begin{frame}{COSO 2017: Integrating with Strategy and Performance}
\begin{itemize}
  \item COSO: joint initiative of five U.S.\ accounting and internal-audit associations; its 1992 \emph{Internal Control} framework became the SOX~404 benchmark.
  \item ERM framework released 2004; updated 2017 as \emph{Enterprise Risk Management: Integrating with Strategy and Performance}.
  \item Core argument: organizations exist to \alert{create value}; ERM gives structure to deciding which risks to accept in pursuit of strategy --- deliberately, consistently, with oversight.
\end{itemize}
\end{frame}

\begin{frame}{Five Components, Twenty Principles}
\begin{enumerate}
  \item \textbf{Governance and Culture} --- the foundation. Board oversight, operating structures, tone at the top. If executives don't actually \emph{use} risk information, the apparatus adds little.
  \item \textbf{Strategy and Objective-Setting} --- \alert{risk appetite} is a strategic choice, not a compliance parameter. Are accepted risks deliberate and consistent with what the board will absorb?
  \item \textbf{Performance} --- the operational work: identify, assess severity (likelihood $\times$ impact), prioritize, respond (accept / avoid / reduce / share), build a \alert{portfolio view}.
  \item \textbf{Review and Revision} --- the risk profile changes as strategy and conditions change.
  \item \textbf{Information, Communication, and Reporting} --- the right information, to the right people, in usable form; communication flows up \emph{and} down.
\end{enumerate}
\medskip
The 20 principles are an interconnected system, not a checklist.
\end{frame}

\begin{frame}{What 2017 Changed}
\begin{itemize}
  \item Earlier treatments framed risk management as protection against downside.
  \item The 2017 update is explicit:
  \begin{itemize}
    \item Value creation \emph{requires} accepting risk.
    \item Risk appetite is a \alert{positive choice} about which uncertainties to embrace.
    \item ERM should help organizations take more \emph{deliberate} risk --- not simply less risk.
  \end{itemize}
  \item That reframing changes how risk officers talk to boards and CEOs about what they do.
\end{itemize}
\end{frame}

\section{The ISO 31000 Standard}

\begin{frame}{ISO 31000 (2009, rev.\ 2018)}
\begin{itemize}
  \item International standard, deliberately generic --- works for a multinational bank, a municipality, a hospital, or a small manufacturer.
  \item Defines risk as \alert{``the effect of uncertainty on objectives''} --- neutral by design; upside and downside both warrant attention.
  \item Three layers:
  \begin{itemize}
    \item \textbf{Principles} --- what effective risk management looks like (integrated, customized, inclusive, dynamic, improvement-oriented).
    \item \textbf{Framework} --- how it embeds in governance, strategy, operations, reporting.
    \item \textbf{Process} --- the step-by-step method: scope and criteria; identify; analyze; evaluate; treat; monitor; communicate throughout.
  \end{itemize}
  \item The \alert{process model} is its most useful contribution --- more explicit about assessment methodology than COSO.
\end{itemize}
\end{frame}

\begin{frame}{COSO vs.\ ISO 31000}
\begin{columns}[T]
\column{0.48\textwidth}
\textbf{COSO ERM 2017}
\begin{itemize}
  \item Origin: U.S.\ accounting profession, internal-control tradition
  \item More prescriptive, governance-oriented
  \item Aligned with U.S.\ boards and SOX-era regulators
  \item Tells you \emph{what must be in place} organizationally
\end{itemize}
\column{0.48\textwidth}
\textbf{ISO 31000:2018}
\begin{itemize}
  \item Origin: international standards body
  \item More flexible, process-focused
  \item Usable across sectors and countries
  \item Tells you \emph{how to run the work} of a risk assessment
\end{itemize}
\end{columns}
\medskip
\begin{block}{In practice: complements, not competitors}
Many organizations structure the program around COSO's five components and run individual risk evaluations with ISO's process. The underlying logic is identical.
\end{block}
\end{frame}

\begin{frame}{A Caveat on Evidence}
\begin{itemize}
  \item Neither framework offers rigorous empirical evidence that following its prescriptions improves performance.
  \item Some supportive findings: Hoyt and Liebenberg (2011) find ERM adoption positively related to Tobin's $q$ for insurers --- but identification is hard and results are not definitive.
  \item Treat the frameworks as \alert{practical guidance}, not validated science. That's a reason for judgment, not for ignoring them.
\end{itemize}
\end{frame}

\section{The Institutional Landscape}

\begin{frame}{Financial Services: Regulation as the Driver}
\begin{itemize}
  \item Banks and insurers were early adopters --- largely because they had no choice.
  \item \textbf{Banks}: Basel~I (1988) minimum capital; Basel~II (2004) three pillars + explicit risk management systems; Basel~III (2010--11) stronger capital and liquidity after 2008. U.S.\ implementation via Fed/OCC/FDIC, plus Dodd-Frank \alert{stress tests}.
  \item \textbf{Insurers}: NAIC \alert{ORSA} (2012) --- forward-looking own-risk and solvency assessments filed annually; EU \alert{Solvency~II} imposes analogous quantitative and governance requirements.
  \item The pattern: risk-calibrated capital, mandatory infrastructure, board-level expectations, stress testing. ERM in financial services is \emph{not optional}.
\end{itemize}
\end{frame}

\begin{frame}{Non-Financial Corporations: Voluntary but Not Casual}
\begin{itemize}
  \item No regulation requires a manufacturer or hospital to run a COSO-aligned program. Pressure converges from several directions instead:
  \begin{itemize}
    \item \textbf{RIMS} --- professional home for risk managers since 1950; the Risk Maturity Model benchmarks programs across seven attributes and five maturity levels.
    \item \textbf{IIA Three Lines Model} (2020) --- operations (first line), risk and compliance functions (second), internal audit as \alert{independent assurance} (third).
    \item \textbf{Rating agencies} --- S\&P, Moody's, Fitch assess ERM in rating methodologies; credible programs support stronger ratings.
    \item \textbf{Investors and proxy advisors} --- ISS and Glass Lewis grade directors partly on board risk oversight.
  \end{itemize}
  \item Boards have fiduciary reasons, auditors professional reasons, raters economic reasons, investors governance reasons --- all pointing the same way.
\end{itemize}
\end{frame}

\section{Implementing ERM}

\begin{frame}{Five Phases --- Not All at Once}
\begin{itemize}
  \item Implementing everything at once produces \alert{compliance theater}: documents that satisfy a checklist but don't change decisions.
\end{itemize}
\medskip
\begin{enumerate}
  \item \textbf{Foundation} --- oversight responsibilities, a coordinating executive, a policy, a common \alert{risk taxonomy} (without shared language, workshops devolve into definitional arguments).
  \item \textbf{Initial risk assessment} --- facilitated workshops; output is the \alert{risk register}. Qualitative is fine at this stage.
  \item \textbf{Ownership and responses} --- a named owner per risk; accept / avoid / reduce / share, guided by board-approved appetite.
  \item \textbf{Integration} --- embed in planning, capital allocation, performance management; quarterly board reporting.
  \item \textbf{Advancing sophistication} --- quantification, stress tests, scenarios, correlations.
\end{enumerate}
\medskip
Phases 1--4 produce most of the value. Don't start at Phase 5.
\end{frame}

\begin{frame}{Structure Matters Less Than Role Clarity}
\begin{itemize}
  \item Centralized CRO vs.\ distributed vs.\ hybrid: less important than clarity about who is responsible for what.
  \item Four roles must be explicit regardless of structure:
  \begin{itemize}
    \item \textbf{Risk owner} --- primary management accountability per risk
    \item \textbf{Risk coordinator} --- CRO/risk function: methodology and oversight
    \item \textbf{Independent assurer} --- internal audit (third line)
    \item \textbf{Governance body} --- board or board risk committee
  \end{itemize}
\end{itemize}
\medskip
\begin{block}{Back to GM}
Ambiguity about any of these roles creates exactly the shared-responsibility gap the ignition switch case illustrates --- in painful detail.
\end{block}
\end{frame}

\begin{frame}{Discussion Questions}
\begin{enumerate}
  \item GM had engineers, lawyers, safety staff, and auditors all touching the ignition switch problem. Which COSO component failed most fundamentally --- and would ISO~31000's process have caught it?
  \item SOX never mentions enterprise risk management. Make the case that it did more for ERM adoption than any framework document --- then make the case against.
  \item A mid-size manufacturer asks whether to build its program on COSO or ISO~31000. What would you ask before answering, and what would you likely recommend?
  \item Why is ERM mandatory in banking and insurance but voluntary nearly everywhere else? Is that distinction defensible on economic grounds?
\end{enumerate}
\end{frame}

\end{document}
